\begin{frame}{Idea 1: Addiction \& Intertemporal Choice}
    \begin{itemize}
        \item \textbf{Research Question}: How does present-bias ($\kappa$) drive the "Relapse Loop" in tobacco and cannabis cessation?
        \item \textbf{Structural Framework}:
            \begin{itemize}
                \item Applying the \alert{Hyperbolic Discounting} model to longitudinal abstinence data.
                \item Identifying the structural barrier: When immediate craving utility $>$ discounted future health utility.
            \end{itemize}
        \item \textbf{Policy Innovation: Dynamic Incentive Design}
            \begin{itemize}
                \item \alert{Front-loaded Rewards}: Shifting cessation subsidies to the "Day 1--7" window where $\kappa$ is most punishing.
                \item \alert{Commitment Contracts}: Testing BDM-based mechanisms for smokers to "pre-commit" to abstinence goals.
            \end{itemize}
    \end{itemize}
\end{frame}

\begin{frame}{Idea 2: Digital Health Equity \& Cognitive Time-Tax}
    \begin{itemize}
        \item \textbf{Research Question}: Do digital health interventions impose a regressive \alert{"Cognitive Time-Tax"} on Medicaid populations?
        \item \textbf{Methodology: Choice Experiment on Nudge Fatigue}
            \begin{itemize}
                \item Attributes: Notification frequency, timing, and framing (loss vs. gain).
                \item Estimating the \alert{MWTA for cognitive load} --- quantifying the point where nudges become "noise."
            \end{itemize}
        \item \textbf{Health Equity Impact}:
            \begin{itemize}
                \item Analyzing if "Nudge Fatigue" is higher in low-income clusters due to baseline administrative burdens.
                \item Optimizing \alert{Incentive-compatible messaging} for high-risk preventive behaviors.
            \end{itemize}
    \end{itemize}
\end{frame}
